\section{Quantitative replacement and completion of the proof}\label{sec:assembly}
We now choose all accuracies in a fixed order. This is necessary
both for a positive logarithmic saving and to avoid feeding a later
density depth into the complexity of the earlier ideal construction.
Write
\[
 \phi=\varphi(L_0),\quad \delta=\frac1{8192\phi},\quad
 \mathcal S_N=\frac{NJ_{\mathcal E}(N)}{AB L^\kappa},\quad L=\log N.
\]
If \(\Gamma_N=0\), the count is zero by its fixed residue condition.
Hence assume \(\Gamma_N>0\); its finitely many possible positive
values have a fixed positive minimum.

\subsection{An explicit discarded-bank exponent}
The identity-class primes have reciprocal mass
\((1/(8\phi)+o(1))\log L\) in the small range and
\((1/(15\phi)+o(1))\log L\) in the big range. Greedily choose
\[
 b_N=\left\lfloor\frac{\log L}{32\phi}\right\rfloor
\]
disjoint banks of each size, each with reciprocal mass \(1+o(1)\).
There is ample mass in both ranges. The overshoot per bank is at
most the largest reciprocal prime, and target-prime deletion is
superlogarithmically small. Thus all masses are \(1+o(1)\), with
total error smaller than every inverse logarithmic power.
By \eqref{eq:success}, each success probability tends to
\(e^{-2}>1/8\). Chernoff's inequality and independence give
\begin{equation}\label{eq:rare-budget}
 \Prob\{\text{fewer than }b_N/16\text{ successes in either range}\}
 \le2e^{-b_N/64}\ll L^{-1/(4096\phi)}=L^{-2\delta}.
\end{equation}
Here \(b_N\ge\log L/(64\phi)\) eventually.
Lemma~\ref{lem:envelope} bounds the discarded correlation by
\(O(\mathcal S_NL^{-2\delta})\), plus exponential and power-smaller
errors. This applies with an indicator or model on either side.
The exact reconstruction \eqref{eq:reconstruction} has total
coefficient mass at most
\[
 3^{2b_N}\prod_g\max(1,V_g^2)\le L
\]
for large \(N\), since the second product is \(1+o(1)\) and
\(2b_N\log3\le(\log3)\log L/(16\phi)\).

\subsection{Moving cutoffs and the reflected target scale}
Set \(\epsilon=L^{-\delta}\). Choose \(0\le\Phi_\epsilon\le1\),
equal to one on \([2\epsilon,1-2\epsilon]\), supported in
\([\epsilon,1-\epsilon]\), with transitions of width \(\epsilon\).
For the harmonic graph comparison put
\(\Psi_\epsilon(t)=t\Phi_\epsilon(t)\) and
\(F_\epsilon(u)=\Psi_\epsilon(e^u)\). Direct differentiation gives
\[
 \|F_\epsilon\|_1\ll1,\qquad
 \|F_\epsilon^{(j)}\|_1\ll_j\epsilon^{1-j}\quad(j\ge1).
\]
At the lower transition the logarithmic derivative cancels the
small transition scale. Fourier integration by parts therefore gives
\begin{equation}\label{eq:mellin-budget}
 \int|\widehat F_\epsilon(\zeta)|(1+|\zeta|)^r\,d\zeta
       \ll_r\epsilon^{-r-1},\qquad
 \int_{|\zeta|>L^3}|\widehat F_\epsilon(\zeta)|\,d\zeta
       \ll_j L^{-(3-\delta)(j-1)}.
\end{equation}
Thus the external phase range stays fixed at \(L^3\); choose the
tail derivative order last.

On a contributing edge of the reflected graph,
\(n<0\), \(n'=n+ND>0\), and
\(t=|n|/(ND)\in[\epsilon,1-\epsilon]\). Exactly
\begin{equation}\label{eq:two-sided-localization}
 \frac{n'}{|n|}=\frac{1-t}{t}
       \in\left[\frac\epsilon{1-\epsilon},
                  \frac{1-\epsilon}{\epsilon}\right].
\end{equation}
For a source dyad \(|n|\asymp X\), insert a smooth target cutoff
supported on \([c\epsilon X,C\epsilon^{-1}X]\) and equal to one
on the contributing range \emph{before} Mellin inversion. Keep it
in every resulting pairing. Removed edges had zero full Mellin
integral originally. Dividing shared labels preserves the ratio
\eqref{eq:two-sided-localization}.

There are \(O(\log L)\) target dyads \(X'\), each with uniformly
bounded rescaled logarithmic derivatives. Their endpoint moments are
\(O(X)\), \(O(X')\): fixing an exact label product and counting its
multiples proves this directly, and summing reciprocal list mass
introduces no exponent depending on the final fixed mark length.
The pairing can therefore cost at most
\(\sqrt{X'/X}\ll\epsilon^{-1/2}\). Together with the unweighted
Mellin norm and dyadic count these costs are at most \(L^{5\delta/2}\).
We reserve \(L^{8\delta}\) for all localization losses.

For the comparison source scale \(Y\asymp ND_1D_2\), use its
actual target dyad \(Y'\), where
\(c\epsilon Y\le Y'\le C\epsilon^{-1}Y\). The physical major
window is \(L^{B_{\rm out}}/Y\), so
\(|\beta|Y'\ll L^{B_{\rm out}+\delta}\). Apply
Proposition~\ref{prop:gr-energy} with that enlarged exponent.
A crude integrated-energy rescaling costs \(\epsilon^{-2}\), still
inside the reserve. At the smallest dyad and selected small prime
\(P\le\exp(CL^{1/4})\),
\begin{equation}\label{eq:smallest-dyad}
 \frac{Y'}{qP}\gg
 e^{cL^{1/3}-CL^{1/4}}L^{-B_{\rm out}-\delta}\gg_A L^A.
\end{equation}
This remains true after every gcd split, since both length and
modulus are divided by the same divisor. Thus neither target
localization nor target valuations spoil the cofactor-energy estimate.

\subsection{The precision order}
Here are sufficient numerical margins. All exponents are fixed
before \(N\to\infty\).
First request normalized graph-pairing error \(L^{-22}\) per source
dyad. At most \(L^2\) source dyads occur, leaving \(L^{-20}\).
Fix the root and ideal accuracies for this request, then the probe
complexity and eventual mark length. The phase range is \(L^3\).

After that finite comparison family is fixed, let \(C_{\rm cmp}\)
dominate its coefficient and dyadic costs, and \(C_w\) its endpoint
pointwise exponent. These constants do not depend on the density
depth. Let \(m_*\) bound free draws per group. Unique factorization
gives a dyadic free-product law
\[
 \nu_i(n)\le L^{K_*}/n,\qquad
 K_*=\frac{\log(m_* !)+m_*\log2}{16\phi}.
\]
Since the law has mass at most one, its squared norm on a dyad of
length \(T_i\) is at most \(L^{K_*}/T_i\). The integer bilinear
bound \eqref{eq:gr-multiplier} consequently permits
\(C_{\rm prod}=2K_*+2\). Choose
\begin{align}
 A_{\min}&>C_{\rm cmp}+2C_w+40+8\delta,\notag\\
 B_{\rm out}&>2A_{\min}+C_{\rm prod}+4.\label{eq:outer-budget}
\end{align}
Dirichlet approximation with denominator ceiling
\(T_1T_2/L^{B_{\rm out}-1}\) yields a multiplier
\(O(L^{-A_{\min}})\) off the grids of denominators \(Nr\),
\(r\le L^{B_{\rm out}}\). The middle terms in
\eqref{eq:gr-multiplier} are superlogarithmically small because
\(T_i\ge\exp(L^{1/3})\).

Next request integrated target-grid energy saving
\[
 T_E>2(C_{\rm cmp}+C_w+40+8\delta).
\]
Use window exponent \(B_{\rm out}+\delta\), and choose its
residue-energy accuracy to pay both the window and the number of
\(r\)'s. Parseval on the other endpoint then gives major contribution
at most \(L^{-T_E/2+C_w+C_{\rm cmp}+8\delta}\le L^{-40}\).
The minor contribution has the same margin by
\eqref{eq:outer-budget}. Exact discrete Parseval, not a supremum
bound, cancels the target-grid factor \(N\).

Finally choose the Hankel depth for the required exceptional-character
centering. If \(C_{\rm low}\) pays record, character and cutoff costs,
take \(3(K+1)/5>T_E+C_{\rm low}+40\).
Choose the derivative order in \eqref{eq:mellin-budget} after its
tail costs. These choices change constants and thresholds, not the
earlier ideal complexity, bounded endpoint moment exponent or phase
range. In particular the precision choices are noncircular.

\begin{proof}[Completion of Theorem~\ref{thm:general}]
After localization, the graph and comparison errors are bounded by
\(O(NL^{-20+8\delta})\). Reconstruction costs at most \(L\), so
the resulting error is
\[
 O(NL^{-19+8\delta})=o(\mathcal S_NL^{-\delta}),
\]
using \(\kappa<2\), \(J_{\mathcal E}\ge1\), and fixed \(A,B\).
The discarded-bank error is \(O(\mathcal S_NL^{-2\delta})\).
Apply the resulting replacement on each side to the identity
\[
 f_U(x)f_V(y)-D_U(x)D_V(y)
   =\xi_U(x)f_V(y)+D_U(x)\xi_V(y).
\]
Proposition~\ref{prop:model-correlation} evaluates the remaining
model correlation. On retained bulk records its logarithm is
\(L+O(\sqrt L+\delta\log L)\), giving relative
\(O(L^{-1/2})\) in leading-density replacement. The sieve uses a
fixed power level; partial summation costs
\(\TV(\Phi_\epsilon)=O(1)\), rather than its largest derivative.
The record and squared-modulus CRT endpoints cost
\(N^{0.45+o(1)}\). The upper endpoint sieve gives
\[
 O(\epsilon\mathcal S_N+N^{0.45+o(1)}).
\]
Its strip length \(N/L^\delta\) exceeds the sieve level. Euler tails
are superlogarithmically small, and
\(\int\Phi_\epsilon=1+O(\epsilon)\). Every error is therefore
\(O(\mathcal S_NL^{-\delta})\), as asserted.

For an interval \(I\), if \(|I|\ll\epsilon\) use the direct upper
sieve. Otherwise clip it to \([\epsilon,1-\epsilon]\) and smooth
its endpoints at width \(\epsilon\). At most four transition strips
of total length \(O(\epsilon N)\) are lost. Total variation remains
\(O(1)\), and the logarithmic derivative norm of a transition at
\(t_0\) is at most \(O_j(t_0^j\epsilon^{1-j})\).
Thus the same budgets apply uniformly in both endpoints of \(I\),
and the cutoff integral is \(|I|+O(\epsilon)\).
Endpoint conventions alter at most two pairs. This proves the
uniform interval assertion and completes the theorem.
\end{proof}
